\documentclass[../main.tex]{subfiles}
\begin{document}

Isosurface extraction is a fundamental problem in scientific visualization, computer graphics, medical imaging, and computational simulation. Given a scalar field defined over a two-dimensional or three-dimensional domain, the objective is to reconstruct a geometric surface corresponding to a specific scalar value, commonly referred to as an isovalue. The extracted surface provides an intuitive representation of volumetric data and enables users to analyze complex spatial structures that are difficult to observe directly from raw numerical values. This chapter presents the research motivation, reviews existing approaches and their limitations, defines the objectives of the thesis, summarizes the main contributions, and outlines the organization of the thesis.

\section{Problem Statement}
\label{sec:dvd}

The rapid development of 3D content creation technologies has increased the demand for efficient methods that convert implicit representations into editable polygonal meshes. Implicit surface representations, such as signed distance fields (SDFs), neural fields, and volumetric scalar fields, have become widely adopted in geometry processing and generative modeling because they can represent complex shapes with high geometric fidelity. However, these representations are not directly compatible with conventional 3D modeling, animation, and rigging workflows, which primarily operate on polygonal meshes.

Recent isosurface extraction methods, including FlexiCubes, have demonstrated the ability to reconstruct high-quality surfaces from implicit fields while maintaining differentiability and geometric accuracy. These characteristics make FlexiCubes particularly suitable for neural shape optimization and 3D reconstruction tasks. Nevertheless, the generated meshes are primarily optimized for geometric reconstruction quality rather than downstream content creation applications.

In practical modeling and animation pipelines, mesh topology plays a critical role. A mesh intended for character rigging or deformation should contain edge loops and polygon distributions that follow the underlying semantic structure of the object. Poor topological organization may lead to undesirable deformations during animation, increased manual retopology effort, and difficulties in subsequent editing operations. Although FlexiCubes produces geometrically accurate surfaces, it does not explicitly consider topological structures required by artists and animators.

As a result, artists often need to perform a separate retopology stage after mesh extraction. This process is time-consuming, requires significant expertise, and becomes a bottleneck when processing large numbers of generated models. The lack of topology-aware isosurface extraction therefore limits the applicability of implicit surface representations in production-oriented 3D workflows.

This thesis addresses the problem of generating topology-aware meshes directly during the isosurface extraction process. Specifically, the research investigates modifications to the FlexiCubes framework that incorporate topological constraints and mesh-structure optimization while preserving geometric fidelity. The goal is to reduce the need for manual retopology and produce meshes that are more suitable for 3D modeling, rigging, and animation applications.


\section{Background and Problems of Research}
\label{sec:giaiphap}

The problem of extracting polygonal surfaces from volumetric or implicit representations has been extensively studied in computer graphics and scientific visualization. Early approaches focused on converting scalar fields into explicit mesh representations that could be rendered efficiently by graphics hardware. Among these methods, Marching Cubes became one of the most widely adopted algorithms due to its simplicity and ability to generate surfaces from volumetric data. Numerous extensions were subsequently proposed to improve surface quality, resolve topological ambiguities, and reduce geometric artifacts.

With the emergence of neural implicit representations, isosurface extraction has gained renewed attention. Modern techniques such as Signed Distance Fields (SDFs), Occupancy Networks, and Neural Radiance Fields represent geometry implicitly rather than as explicit polygonal meshes. These representations enable compact storage and high-fidelity reconstruction of complex shapes. However, practical applications in computer graphics still require explicit mesh representations for rendering, editing, simulation, and animation. Consequently, efficient and accurate extraction of polygonal surfaces remains an essential component of modern geometry processing pipelines.

Recent advances in generative artificial intelligence have further accelerated the development of automatic 3D content creation systems. Models such as TRELLIS generate sparse volumetric latent representations from text or image inputs, which are subsequently decoded into polygonal meshes using differentiable isosurface extraction algorithms such as FlexiCubes. This pipeline enables efficient generation of high-quality 3D assets while maintaining compatibility with gradient-based optimization.

\begin{figure}[H]
	\centering
	\includegraphics[width=\textwidth]{Figure/Figure1-Pipeline.png}
	\caption{Overview of the proposed topology-aware mesh generation pipeline. }
	\label{fig:overview_pipeline}
\end{figure}

As illustrated in Figure~\ref{fig:overview_pipeline}, existing generation pipelines are primarily designed to maximize geometric reconstruction quality. Although the generated meshes accurately reproduce the target surface, the extracted topology is largely determined by the local geometry and the underlying extraction algorithm rather than by desirable modeling characteristics.

FlexiCubes is a representative differentiable isosurface extraction method. By introducing flexible vertex placement and learnable optimization mechanisms, it significantly improves reconstruction quality while preserving compatibility with gradient-based learning systems. Experimental results demonstrate that the method produces meshes with higher geometric fidelity than conventional extraction techniques and has therefore been adopted in several modern 3D generation frameworks.

Despite these advantages, existing isosurface extraction methods, including FlexiCubes, primarily optimize geometric reconstruction accuracy. The generated meshes are commonly evaluated using metrics such as surface approximation error, normal consistency, and reconstruction quality. In contrast, considerably less attention is given to the topological structure of the resulting mesh from the perspective of digital content creation.

In practical 3D modeling and animation workflows, mesh topology is often as important as geometric accuracy. A well-structured mesh should contain edge flows that follow the semantic features of the object and support deformation during animation. Character models, for example, require organized edge loops around joints and facial regions to achieve natural deformations during rigging and skinning. Similarly, editable assets benefit from mesh structures that simplify sculpting, subdivision, texture mapping, and geometric editing. Meshes generated solely for geometric accuracy frequently contain irregular vertex distributions and edge connectivity that are difficult to use in these downstream tasks.

As a consequence, assets generated from implicit representations often require a separate retopology stage before they can be integrated into production pipelines. Retopology reconstructs a more suitable mesh structure while preserving the original geometry. Although effective, this process introduces additional computational cost and often requires substantial manual effort from artists. The need for post-processing reduces the overall efficiency of workflows based on automatic geometry generation.

The limitations discussed above reveal a gap between existing isosurface extraction methods and the requirements of practical 3D content creation. While FlexiCubes successfully addresses challenges related to differentiable surface reconstruction and geometric fidelity, it does not explicitly optimize mesh topology. This observation motivates the research presented in this thesis, which investigates topology-aware modifications to the FlexiCubes framework by incorporating directional priors learned from high-quality reference meshes. The objective is to preserve the geometric advantages of differentiable isosurface extraction while generating meshes with more consistent edge flow and improved suitability for modeling, subdivision, rigging, and animation.


\section{Research Objectives and Conceptual Framework}

The objective of this thesis is to investigate the limitations of current isosurface extraction methods in the context of 3D content creation and to develop an extension of the FlexiCubes framework that generates topology-aware meshes suitable for modeling and rigging applications. While existing methods focus primarily on reconstructing geometry with high accuracy, the proposed research aims to incorporate topological considerations into the mesh generation process without significantly compromising reconstruction quality.

To achieve this objective, the thesis first studies the theoretical foundations of implicit surface representations, mesh topology, and isosurface extraction algorithms. Existing approaches are then analyzed to identify the factors that influence mesh structure, edge flow, and vertex distribution. Particular attention is given to the FlexiCubes algorithm due to its differentiable formulation and its ability to produce high-quality surface reconstructions from implicit fields.

Based on this analysis, the thesis proposes a modified FlexiCubes framework that integrates topology-aware constraints into the extraction process. Instead of optimizing solely for geometric reconstruction accuracy, the proposed approach considers additional criteria related to mesh structure and topology quality. These criteria aim to generate meshes that are more compatible with common operations in 3D modeling, subdivision, skinning, and animation workflows.

The proposed framework consists of three main stages. In the first stage, an implicit surface representation is obtained from the input data. In the second stage, a modified FlexiCubes extraction process is performed to reconstruct the target surface while incorporating topology-aware optimization criteria. In the final stage, the generated mesh is refined through topology-preserving operations to improve its suitability for downstream applications.

The effectiveness of the proposed method is evaluated through experiments on representative 3D models. The resulting meshes are compared with those generated by the original FlexiCubes algorithm using both geometric and topological metrics. The evaluation focuses on reconstruction accuracy, mesh quality, topological regularity, and suitability for modeling and rigging tasks. Through these experiments, the thesis seeks to determine whether topology-aware modifications can reduce the need for manual retopology while preserving the advantages of differentiable isosurface extraction.


\section{Contributions}

This thesis provides the following main contributions:

\begin{enumerate}
	\item A comprehensive review and analysis of representative isosurface extraction algorithms, including their theoretical foundations, advantages, and limitations.
	
	\item An improved isosurface extraction framework designed to enhance geometric accuracy and topological consistency while maintaining computational efficiency.
	
	\item An experimental evaluation on volumetric datasets that demonstrates the effectiveness of the proposed approach through comparisons with existing methods.
\end{enumerate}

\section{Organization of the Thesis}

The remainder of this thesis is organized as follows.

Chapter~2 presents the theoretical background and related work that form the foundation of this research. The chapter first defines the scope of the study and reviews representative isosurface extraction algorithms, including Marching Cubes, Dual Contouring, Deep Marching Cubes, DMTet, and FlexiCubes, highlighting their advantages and limitations. It then introduces the datasets used in this thesis, namely ShapeNetCore and TRELLIS-generated geometry, followed by the background knowledge on Transformer architectures, implicit surface representations, isosurface extraction, and mesh topology. The chapter concludes by identifying the research gap that motivates the proposed topology-aware isosurface extraction framework.

Chapter~3 describes the proposed methodology for topology-aware isosurface extraction. The chapter first presents the extraction of directional information from high-quality meshes in the ShapeNetCore dataset and explains how local edge-flow information is transformed into a continuous directional field. Next, it introduces a Transformer-based neural network trained to predict surface normals and edge directions from point clouds. Finally, the chapter presents the integration of the pretrained model into the FlexiCubes optimization framework, where the signed distance field, vertex deformation, and adaptive cube parameters are jointly optimized to generate meshes whose edge flow better matches the learned topology priors while preserving geometric fidelity.

Chapter~4 presents the experimental setup and evaluation methodology used to assess the proposed approach. The chapter describes the evaluation metrics, datasets, implementation details, and experimental procedures. Quantitative comparisons between the original FlexiCubes method and the proposed topology-aware framework are conducted using geometric reconstruction metrics, including Root Mean Square (RMS) error, Hausdorff distance, and surface normal deviation. The experimental results are then analyzed to evaluate the effectiveness of the proposed optimization strategy and discuss its advantages and limitations.

Finally, Chapter~5 summarizes the main contributions of the thesis and discusses the overall findings of the research. The chapter highlights the achievements and remaining limitations of the proposed method and concludes by suggesting several directions for future work, including extending the approach to more complex object categories, improving topology guidance using more flexible grid representations, and integrating topology-aware optimization directly into generative 3D models.


\end{document}